\documentclass[solutions]{worksheet}
\usepackage[english]{babel}
\usepackage{amsmath}

\coursename{Number Theory}
\semester{Fall 2026}
\weeknumber{4}
\worksheetnumbertotal{10}
\worksheetnumber{4}

\begin{document}

\begin{exercise}

  Is $91$ prime? Justify your answer.

\end{exercise}

\begin{solution}

  No: $91 = 7 \times 13$, so it is not prime.

\end{solution}

\begin{exercise}

  This exercise has multiple sub-questions about modular arithmetic.

  \begin{questions}

    \question Compute $2^{10} \bmod 11$.

    \question State Fermat's little theorem and explain why your answer to
    the previous question is consistent with it.

  \end{questions}

  These ideas are the foundation of several primality tests used in modern
  cryptography.

  \begin{questions}

    \question Does the converse of Fermat's little theorem hold, i.e.\ if
    $a^{p-1} \equiv 1 \pmod p$ for some $a$, must $p$ be prime?

  \end{questions}

\end{exercise}

\begin{solution}

  Fermat's little theorem is a cheap first filter for primality: it never
  wrongly rejects a prime, but it can be fooled by certain composite
  numbers (Carmichael numbers), which is why the converse fails and why
  real primality tests use several independent bases.

  \begin{answers}

    \answer By Fermat's little theorem, $2^{10} \equiv 1 \pmod{11}$ since
      $11$ is prime and $\gcd(2, 11) = 1$.

    \answer Fermat's little theorem states that if $p$ is prime and
      $\gcd(a, p) = 1$, then $a^{p-1} \equiv 1 \pmod p$. Since $11$ is
      prime, this predicts exactly the result computed above.

    \answer No. A composite number $p$ for which $a^{p-1} \equiv 1
      \pmod p$ holds for some (or even every) $a$ coprime to $p$ is called
      a pseudoprime (or, if it holds for every valid $a$, a Carmichael
      number) -- for example $341 = 11 \times 31$ satisfies
      $2^{340} \equiv 1 \pmod{341}$ despite being composite.

  \end{answers}

\end{solution}

\end{document}
